% ============================================================
% kbs_disc.tex — KBS journal: Discussion and limitations (Section 9).
% Section label: sec:disc.
% Bibliography: kbs_ke_references.bib + paper/references.bib.
% ============================================================

\section{Discussion and Limitations}
\label{sec:disc}

\subsection{What the paper establishes}

Read as a knowledge-engineering study, this paper establishes three
claims. First, LLM-acquired subject knowledge is a \emph{viable
acquisition channel}: validated against professional gold under a
multi-level rubric that respects legitimate inter-cataloger variation, it
operates near the agreement ceiling the gold itself imposes, and its errors
are rarely fabrication (Section~\ref{sec:valid}) \TODO{E2: report panel
ceiling}. Second, its utility is \emph{representation-dependent}: a
knowledge-based retrieval system consumes the thesaurus representation of
the acquired knowledge and not the hierarchical one, and a knowledge
manager targeting topical discovery should therefore acquire subject
headings, not classification (Section~\ref{sec:util}). Third, its value is
\emph{completeness-dependent}: the gain from acquisition is largest where
the knowledge base is sparsest and never reaches zero, which locates the
rational deployment in the gap-filling regime (Sections~\ref{sec:util}
and~\ref{sec:synth}).

The methodological contribution is separable from these findings: the
paper completes the knowledge-engineering cycle for machine-acquired
knowledge---acquire, represent, validate against consensus, measure in the
consuming system---and reports each stage against the ceiling the domain
imposes, with every number traceable to released artifacts
(Sections~\ref{sec:data}--\ref{sec:cost}).

\subsection{Generalization}

The domain is bibliographic subject knowledge, but the object of study is
the acquisition channel. Three results should transfer to other knowledge
domains, and one should not. \emph{Representation dependence}---the finding
that downstream utility is carried by an open, pre-coordinated vocabulary
and not by a closed hierarchy---is a claim about the structure of
knowledge representations and the retrieval task, not about libraries;
we would expect an analogous contrast wherever a consuming system matches
queries against one vocabulary type. \emph{Completeness dependence}---value
concentrated where coverage is low---should transfer to any knowledge base
with uneven coverage. The \emph{validation methodology} (multi-level
agreement read against an inter-source ceiling) transfers directly. What
should not transfer is the specific gain magnitude, which is collection- and
task-specific; we report it within the studied collections and make no
claim about collections unlike them.

\subsection{Limitations}

\textbf{Query realism.} Topical questions in the discovery pool are
paraphrases of each record's own gold subject heading. Pooled graded
judging and human calibration (Section~\ref{sec:util}) fix the scoring
bias, not the query distribution: topical numbers remain an upper bound on
what unconstrained patron queries would yield. The authority-sourced query
design---sampling headings independently of the corpus and judging
relevance to them---is the planned fix \TODO{run authority-sourced query
set}.

\textbf{Corpus and language.} The collection is canonical English
public-domain literature. The memorization probe bounds recall effects
across well-known works but says nothing about genuinely obscure or
contemporary material; a corpus with such material is required to extend
the acquisition claim beyond the studied regime \TODO{E1}.

\textbf{Gold and agreement.} Gold subject knowledge is professional but
single-source at benchmark scale; the HathiTrust cross-check quantifies
inter-source disagreement, and the panel-validated gold and measured
agreement ceiling \TODO{E2} convert the single-source figures into
consensus-relative accuracy. Classification gold aggregates editions and is
scored as any-gold-match.

\textbf{Models and conditioning.} The primary acquisition runs use a single
LLM conditioned on sparse records. Multi-model replication and
full-text-conditioned acquisition \TODO{E3} bound the model and conditioning
dependence of the validity results; retrieval results are
model-independent by construction (they are scored against record
identifiers with no LLM in the loop).

\textbf{Cost estimates.} The cost side rests on a small expert panel and is
reported per-cataloger with sensitivity analysis \TODO{E4}; the deployment
analysis should be read as a pilot decision model, not a precise
cost schedule.

\textbf{Representation coverage of the study.} We study one thesaurus
(LCSH) and one taxonomy (DDC) as instances of the two representation
classes. The mechanism argument---open pre-coordinated strings versus
closed hierarchies---predicts the contrast generalizes across vocabularies,
but only the two instances are measured here.

\subsection{Ethics and responsible AI}

All source data are open and public-domain; no patron or circulation data
is used. Generated metadata is evaluated for research purposes only and
never written back to production catalogs. The generative-AI disclosure
(Section~\ref{sec:method:benchmark}) states every role the LLM played,
including as judge and as question polisher, and all generated content was
programmatically validated. The cataloger-panel study follows an explicit
protocol with consent and de-identification \TODO{E2: attach the IRB/consent
note}. We report regressions as prominently as gains---the
bibliographic-fact deficit and the known-item trade are in the results, not
the appendix---because the deployment decision this paper supports is only
as trustworthy as its honest cost accounting.

\subsection{Future work}

Beyond the extensions already flagged (authority-sourced queries,
multi-model and full-text-conditioned acquisition, panel-validated gold and
cost), three directions follow directly. First, the representation-dependence
result should be tested across other vocabulary pairs (other thesauri,
other taxonomies, and ontologies proper) to confirm the mechanism.
Second, the decision model in Section~\ref{sec:synth} should be calibrated
against real collections with measured coverage distributions rather than
the simulated coverage sweep. Third, the acquisition component should be
studied in a fully assisted workflow---draft, verify, and authority-gate in
one loop---where the interaction of the three stages is itself the
treatment.
